% Grievance Policy.
%
% NOT derived from a .docx. POSH tells an intern where to take a sexual
% harassment complaint and nothing told them where to take any other one --
% about their supervisor, their work, their treatment, or this documentation.
% An unpaid intern is the least likely person in an organisation to complain
% and the most in need of a route that does not run through the person they
% are complaining about.
%
% WORDING IS A DRAFT FOR REVIEW BY QUALIFIED COUNSEL.
\documentclass{paridyum}
\begin{document}

\paridyumletterhead
\paridyumdateref{\dkDocDate}{\dkDocRef}
\paridyumtitle{Grievance Policy}

\textit{This policy explains how to raise a concern about your work, your
treatment, or a decision affecting you, and what Paridyum will do about it. It
applies to everyone working at or with Paridyum, including interns.}

\section*{1. What This Covers}
Any concern about your engagement with Paridyum, including:
\begin{itemize}[leftmargin=*, noitemsep]
    \item how you have been treated by anyone at Paridyum;
    \item the work you have been assigned, or the hours expected of you;
    \item a decision about your internship, including its ending;
    \item the documents you were asked to sign, or how your personal data has
          been handled; and
    \item anything you believe is unsafe, dishonest, or unlawful.
\end{itemize}

\textbf{Sexual harassment is not dealt with here.} It has its own route, with
its own protections, under the Policy on Prevention of Sexual Harassment. You
may use that route directly without raising anything under this policy first.

\section*{2. Where to Take It}
\begin{infobox}
\textbf{Ordinarily, your supervisor} --- \dkPEngagementSupervisor{} --- if the
concern is about your work.

\vspace{0.5em}
\textbf{If the concern is about your supervisor}, or you would rather not go to
them, write directly to \textbf{\dkOrgEmail}. It reaches the officers and not
your supervisor.

\vspace{0.5em}
\textbf{If your concern is about an officer of Paridyum}, write to the other
officer named on this document. You are not required to raise a concern with
the person it is about.
\end{infobox}

\section*{3. What Paridyum Will Do}
\begin{itemize}[leftmargin=*, noitemsep]
    \item \textbf{Acknowledge} your concern within \dkOrgGrievanceAcknowledgeDays{}
          working days, in writing.
    \item \textbf{Look into it} --- talking to you, and to anyone else who needs
          to be asked.
    \item \textbf{Tell you the outcome} in writing within
          \dkOrgGrievanceOutcomeDays{} working days, or tell you why more time
          is needed and when to expect an answer.
    \item \textbf{Tell you what changes}, if anything does.
\end{itemize}

\section*{4. If You Are Not Satisfied}
Say so, in writing, to the officers. They will review the outcome and respond.
Nothing in this policy prevents you from raising a matter with an external
authority, and Paridyum will not treat doing so as a breach of any obligation
you owe it.

\section*{5. No Retaliation}
No one will be disadvantaged for raising a concern in good faith, for
supporting someone else's, or for giving information during one. This applies
fully to interns: an internship will not be shortened, altered, or ended, and a
certificate or reference will not be withheld or weakened, because someone
raised a concern. Retaliation is itself a breach of the Code of Conduct.

This paragraph is the reason the policy exists. An unpaid intern on a fixed
term has the least security of anyone here, and a route to complain is worth
nothing if using it is a risk.

\section*{6. Confidentiality}
A concern is discussed only with those who need to know in order to deal with
it. Where you ask for something to be kept confidential, Paridyum will do so as
far as it can, and will tell you if it cannot.

\section*{7. Anonymous Concerns}
You may raise a concern without giving your name, by writing to
\dkOrgEmail. Paridyum will look into it as far as it can, but it may not be
able to reach a conclusion or tell you the outcome.

\section*{8. Acknowledgement}
By signing below, the Intern confirms that they have received this policy and
know where to take a concern.

\vspace{1.2em}
\paridyumofficers{\dkSigCeoName}{\dkSigCeoDesignation}{\dkSigCtoName}{\dkSigCtoDesignation}

\vspace{1.2em}
\paridyumcountersign{Intern}{\dkPIdentityLegalName}

\end{document}
